\documentclass[11pt]{article}

\usepackage[margin=1in]{geometry}
\usepackage[T1]{fontenc}
\usepackage{lmodern}
\usepackage{amsmath}
\usepackage{booktabs}
\usepackage[table]{xcolor}
\usepackage{cleveref}

\title{Colored Tabular Rules under \texttt{colortbl}:\\A Fidelity Stress Test}
\author{Test Corpus}
\date{}

\begin{document}
\maketitle

\section{Overview}

This fixture exercises the interaction between the \texttt{booktabs} package and
the \texttt{colortbl} mechanism loaded through \verb|\usepackage[table]{xcolor}|.
Color directives in a tabular environment are notoriously difficult to reproduce
in a word-processing target because they operate at four distinct granularities:
whole rows, individual cells, entire columns, and inline text spans. A converter
that flattens a colored table into a plain grid loses information that a careful
reader relies on, so each granularity is represented separately in
\cref{tab:budget} below.

The table reports a synthetic departmental budget across three fiscal years. Row
shading alternates to aid horizontal tracking; one cell is highlighted to draw
attention to an out-of-band figure; one column carries a faint background tint
applied through the column specification; and a single entry uses colored text
to flag a provisional value. None of these effects changes the underlying
numbers, which means a faithful conversion should preserve both the data and the
visual encoding that distinguishes the rows.

\begin{table}[htbp]
\centering
\caption{Synthetic departmental budget with row, cell, column, and text coloring.}
\label{tab:budget}
\begin{tabular}{l >{\columncolor{blue!8}}r r r}
\toprule
Category & FY2023 & FY2024 & FY2025 \\
\midrule
\rowcolor{gray!20}
Personnel      & 1{,}240 & 1{,}310 & 1{,}388 \\
Equipment      &   305 &   290 &   332 \\
\rowcolor{gray!20}
Travel         &    88 &    95 & \cellcolor{yellow} 412 \\
Software       &   146 &   158 &   171 \\
\rowcolor{gray!20}
Contingency    &    50 &    50 & \textcolor{red}{75} \\
\midrule
Total          & 1{,}829 & 1{,}903 & 2{,}378 \\
\bottomrule
\end{tabular}
\end{table}

\section{Reading the encoding}

The alternating \verb|\rowcolor{gray!20}| shading marks the Personnel, Travel,
and Contingency rows, leaving the intervening rows on the default background. The
\verb|\cellcolor{yellow}| applied to the FY2025 travel figure signals that the
\$412 outlay is anomalous relative to the prior two years, when travel hovered
below \$100. Because the second column is declared with
\verb|>{\columncolor{blue!8}}r|, every FY2023 figure sits on a faint blue tint
regardless of any row shading, which lets a reader follow that column down the
page. Finally, the provisional contingency figure for FY2025 is set in
\verb|\textcolor{red}{75}| so that it reads as an estimate rather than a
committed amount.

Taken together, the four mechanisms in \cref{tab:budget} cover the practical
range of color usage in scholarly tables. A conversion that drops the column
tint, collapses the alternating rows to a single fill, or silently discards the
red text would degrade the information content even while leaving the numerals
intact, which is precisely the failure mode this fixture is designed to surface.

\end{document}
